\chapter{Background}

\section{Real-Time Systems}
A real-time system is defined as a system which requires correct logical operation with compeltion a configured deadline. In contrsast to a non-real-time system which correct operation is desired but lacks configured deadlines. \cite{kaviRealTimeSystems2009}

These real-time systems can be categorised into hard and soft. Hard real-time systems contain strict timing constraints, where a response must be calculated in a specified time-frame and a missed deadline could result in a catestrophic outcome. Soft real-time systems also have timing constraints, but a missed deadline should not result in a catestrophic outcome.

Typical hard real-time systems include aircraft or robotic control computers and soft real-time systems include media players or gaming computers. Soft real-time systems can be used in control applications, but the soft timing must be considered in the system design.

\section{Control Systems}
A control system
Differetial equation

\section{Operating Systems}
An Operating System (OS) is a piece of software that controls the execution of application programs on a system. Usually this involves managing resource allocations, such as a processor execution time, memory usage and generic inputs/outputs. \cite{stallings2018operating}

In order to develop a real-time applications with an OS, a RTOS (Real-Time Operating System) is usually required. This is to ensure the time deadlines of an application are inforced.

Hard real time operating system are designed to run applications with absolute timing
Worst-Case Execution Time
Assigned a realtive deadline, all tasks must complete in their deadline.
Modelled mathematically to ensure that all tasks run to completion

Soft Real Time

Define the maxiumum delay after a deadline


Scheduling
Memory

\section{Mixed Criticality Systems}



\section{Middleware}


\section{Software Updates}